\documentclass[11pt]{article}
\usepackage[a4paper,margin=27mm]{geometry}
\usepackage{amsmath,amssymb,amsthm,xcolor,booktabs,array}
\newcommand{\PROVED}{\textcolor{green!40!black}{\textbf{PROVED}}}
\newcommand{\OPEN}{\textcolor{red!70!black}{\textbf{OPEN}}}
\newcommand{\AUDIT}{\textcolor{orange!70!black}{\textbf{AUDIT}}}
\newtheorem{theorem}{Theorem}
\newtheorem{corollary}{Corollary}
\title{A positive Witt--Fock reservoir with strict unit budget (v110)}
\date{13 August 2026}

\begin{document}
\maketitle

\section{The two-orientation word Gram}

For $N\ge2$ and $k\ge1$ let
\[
 \mathscr K_{N,k}=\ell^2\{1,\ldots,N\},
 \qquad
 v_{N,k}(m)=\frac{\Lambda_k(m)}{\sqrt m(\log N)^k}.
\]
If $\mathcal C_N=\sum_{n\le N}\Lambda(n)n^{-1/2}S_n$ is the finite
logarithmic connection of G39, then Dirichlet convolution gives the exact
Krylov identity
\begin{equation}
                 v_{N,k}=(\log N)^{-k}\mathcal C_N^k\delta_1. \tag{0}
\end{equation}
Thus the reservoir below is generated by the already constructed G39
connection; no spectral zero or desired Schur sign enters its definition.
Let $P_N$ be the coordinate projection onto the divisors of $N$, and let
$R_N$ be the selfadjoint unitary on $P_N\mathscr K_{N,k}$ defined by
$R_Ne_m=e_{N/m}$.  In the doubled space
$\mathscr G_{N,k}=\mathscr K_{N,k}\oplus\mathscr K_{N,k}$ put
\begin{equation}
 x_{N,k}=(v_{N,k},0),
 \qquad
 y_{N,k}=(R_NP_Nv_{N,k},(I-P_N)v_{N,k}).                      \tag{1}
\end{equation}

\begin{theorem}[Exact positive realization of $V_{N,k}$ and $H_{N,k}$ ---
\PROVED]
The two columns in (1) have Gram matrix
\begin{equation}
 \begin{pmatrix}
  \langle x_{N,k},x_{N,k}\rangle&
  \langle x_{N,k},y_{N,k}\rangle\\
  \langle y_{N,k},x_{N,k}\rangle&
  \langle y_{N,k},y_{N,k}\rangle
 \end{pmatrix}
 =\frac1{(\log N)^{2k}}
 \begin{pmatrix}V_{N,k}&H_{N,k}\\H_{N,k}&V_{N,k}\end{pmatrix}. \tag{2}
\end{equation}
Consequently the symmetric and antisymmetric innovations
\begin{equation}
 w_{N,k}^{\pm}=\frac{x_{N,k}\pm y_{N,k}}{\sqrt2}              \tag{3}
\end{equation}
are orthogonal and satisfy
\begin{equation}
 \|w_{N,k}^{\pm}\|^2
 =\frac{V_{N,k}\pm H_{N,k}}{(\log N)^{2k}}.                  \tag{4}
\end{equation}
\end{theorem}

\begin{proof}
The first norm in (2) is the definition of $V_{N,k}$.  Orthogonality of
$P_N$ and $I-P_N$, together with unitarity of $R_N$ on the divisor
coordinates, gives $\|y_{N,k}\|=\|v_{N,k}\|$.  Finally
\begin{align*}
 \langle x_{N,k},y_{N,k}\rangle
 &=\sum_{m\mid N}
 \frac{\Lambda_k(m)\Lambda_k(N/m)}{\sqrt N(\log N)^{2k}}\\
 &=\frac{\Lambda_{2k}(N)}{\sqrt N(\log N)^{2k}},
\end{align*}
by Dirichlet convolution.  This proves (2), and diagonalizing the displayed
$2\times2$ Gram proves (3)--(4).
\end{proof}

\section{All depths}

Define the positive word reservoir and its symmetric source vector by
\begin{equation}
 \mathscr F_N=\bigoplus_{k\ge1}\mathscr G_{N,k},
 \qquad
 W_N^+=\bigoplus_{k\ge1}w_{N,k}^+.                           \tag{5}
\end{equation}
Only finitely many summands are nonzero for each $N$.

\begin{corollary}[Strict asymptotic reservoir norm --- \PROVED]
The vector $W_N^+$ is source-defined from the two arithmetic orientations
and
\begin{equation}
 \|W_N^+\|^2
 =\sum_{k\ge1}\frac{V_{N,k}+H_{N,k}}{(\log N)^{2k}}
 \longrightarrow
 \frac{\sqrt\pi}{2}e^{1/4}\operatorname{erf}\!\left(\frac12\right)
 =0.5922965364\ldots<1.                                      \tag{6}
\end{equation}
In particular there are $N_0$ and $\eta>0$ such that
\begin{equation}
                         \|W_N^+\|^2\le1-\eta
 \qquad(N\ge N_0).                                           \tag{7}
\end{equation}
\end{corollary}

\begin{proof}
Equation (4), summed over $k$, gives the first equality.  The strict limit
and an explicit limiting margin $0.4077034635\ldots$ are Theorem 2 of v109.
Choose any $\eta$ smaller than that margin and then take $N_0$ large enough.
\end{proof}

The construction is stronger than a scalar majorant: (2) retains both
orientations and their exact convolutional cross term.  It is also
independent of the desired sign of the physical operator.

\section{The remaining synthesis map}

Let $\mathscr E_N$ denote the newborn physical boundary cell and let
\begin{equation}
 \mathcal R_Ne=D_0^{\dagger/2}(D_0h_e-q_e)                    \tag{8}
\end{equation}
be its normalized Schur return.  The object which would complete the bridge
is now concrete: a source-defined linear synthesis
\begin{equation}
 \Sigma_N:\mathscr F_N\widehat\otimes\mathscr E_N
       \longrightarrow\overline{\operatorname{Ran}D_0^{1/2}} \tag{9}
\end{equation}
such that
\begin{equation}
 \mathcal R_Ne=\Sigma_N(W_N^+\otimes e),
 \qquad \|\Sigma_N\|\le1+o(1).                              \tag{10}
\end{equation}
If (10) is obtained from the already fixed G39 divergence and the
Gamma--Tate terminal normalization, then (6) gives
\begin{equation}
 \|\mathcal R_Ne\|^2\le(0.5922965364\ldots+o(1))\|e\|^2,     \tag{11}
\end{equation}
which is a strict large-cell Schur margin.  The remaining finitely many
cells can then be treated by the certified Feshbach method of row (d).

Defining $\Sigma_N$ from (8) and declaring it contractive would be
circular.  A valid proof must give its action first on the divisor-incidence
generators of G39 and verify the norm using the Gram (2).  Thus (9)--(10)
is a typed operator construction, not another scalar reformulation of
condition $G$.

\begin{center}
\begin{tabular}{>{\raggedright\arraybackslash}p{.55\linewidth}
                >{\centering\arraybackslash}p{.15\linewidth}
                >{\raggedright\arraybackslash}p{.20\linewidth}}
\toprule
Statement & Status & Location \\
\midrule
Positive two-orientation Witt word Gram & \PROVED & (1)--(4) \\
Identification with the G39 Krylov orbit & \PROVED & (0) \\
Strict all-depth unit budget & \PROVED & (5)--(7), v109 \\
Physical synthesis $\Sigma_N$ with (10) & \OPEN & G39/Gamma--Tate interface \\
Large-cell G40 margin & \OPEN & follows from (10)--(11) \\
Condition $G$, row (d) & \OPEN & large cells plus finite certification \\
\bottomrule
\end{tabular}
\end{center}

\end{document}
